% ************************** Thesis Abstract *****************************
\begin{abstract}
% TODO: this abstract still describes the proposal (``anticipated outcome). Rewrite it after the results are final.
The transport sector is the largest contributor to environmental issues related to climate change and air pollution. Ridesharing has shown promise to promote sustainable transportation, but this service is not widely adopted in Cebu, Philippines. This thesis project details the development and deployment of a ridesharing mobile application tailored for the student community of the University of San Carlos (USC). The project aimed to establish an accessible, safe, and sustainable transportation service for student commuters, addressing the environmental impact of conventional road transport. An in-depth literature review focusing on existing ridesharing algorithms, student transportation behaviors, and the value of ridesharing within the Southeast Asian context was synthesized for key insights to guide development. The methodology encompassed the design of a conceptual framework, specification of core matching and routing algorithms, and definition of the system flow for the mobile application. The result was a functional ridesharing platform designed to meet the specific transportation needs of USC students while promoting environmentally conscious commuting practices.

\medskip
\noindent\textbf{Keywords:} ridesharing, application, transportation, ride-matching, students
\end{abstract}
